\section{Gamma convolution descent of the original arithmetic source}
\label{sec:gamma-convolution-descent}

\textbf{Purpose and result.}
The Toda calculation expresses the original arithmetic control through
source and relation Gram determinants. The present calculation makes their
input explicit at every tensor degree. It keeps the gamma reference measure
already used for the theta source, multiplies by the actual arithmetic
amplitude, and carries both through the literal sum of the spectral variables.
The sum projection comes with its complete relative-fibre complement.
An exact one-factor coefficient transform supplies every finite source and
relation moment; a gamma majorant supplies an explicit tail bound.

The full proof body of \emph{An exact gamma pushforward for the arithmetic
Toda control} (13 September 2026) follows. Its equations are prefixed GD
and its headings are subordinate here. Its inspection statements describe
the source's pinned historical GitHub read, not a later CI or publication
status. The added final subsection proves the direct analytic map from
the coefficient series to the arithmetic moment function used by Toda.
The arithmetic multiplier, its complex phase, the total mass, the original
cyclic relation, and the supported zero all remain in the displayed maps.

\textbf{Degree domains.}
For a nonzero cyclic module of dimension \(q\), the quotient metrics and
representative maps in GD.31--33 are defined for \(N\geq q-1\).
The adjacent-volume ratios and control expression in GD.34--35 require
\(N\geq q\). At \(N=q-1\) the separate first-degree control formula in
the preceding Toda chapter applies; no \(V_{q-2}^{-1}\) is introduced.
For \(h=1\) the finite quotient is zero-dimensional, while the analytic
seed and its mass remain nonzero. No growing-degree vanishing estimate
or conclusion about all zeta zeros is asserted.

\begingroup
\setlength{\emergencystretch}{3em}
\providecommand{\tightlist}{\setlength{\itemsep}{0pt}\setlength{\parskip}{0pt}}

\subsection[\hspace{0.4em}Scope and the inspected collaboration
state]{Scope and the inspected collaboration
state}\label{gammad-scope-and-the-inspected-collaboration-state}

This is a mathematical continuation of \textbf{Tau Toda Volume Control}
and \textbf{Tau Exterior Trace Amplification}. It retains their
arithmetic source, quotient, scaling action, full packet multiplicities,
and canonical minimum-norm metric. The new calculation links those
objects to the gamma reference that is already in the repository. The
gamma orthogonal polynomials and their general convolution identities
are classical. The contribution here is the explicit arithmetic
coefficient transform, its source and quotient maps, its quantitative
bounds, and the resulting exact substitution into the existing
Toda/exterior estimate.

The GitHub read was pinned to main
\path{fa4be32f87a9a90f3c02a07f7790dea95a1e7b0f}. The separate
formalization was read at PR \#22 head
\path{811210d24b80813a08972ca23f919db015383137}. At that inspected
revision the PR explicitly said strict verification was still in
progress. No successful certificate for that revision is inferred here.
Its invariant-filtration trace budget is retained as the other session's
result, not presented as a new deduction of this note.

The main-branch gamma source read was
\path{workbenches/splitzero-tandem/continuations/20260912-stieltjes/tex/theta_gamma_reference.tex},
equations TG.2--TG.14. It already constructs the exact one-factor gamma
measure, its constants, its monic polynomials and its arithmetic
multiplier. We extend that specified reference through the sum
pushforward at every tensor degree.

There is no new RH proof or arithmetic asymptotic assertion here.
Theorems below have written proofs. The checker supplies finite exact
regression evidence, not analytic interval certification and not Lean
execution. No inherited file or other-session branch is modified.

\subsection[\hspace{0.4em}The original arithmetic source and its split
maps]{The original arithmetic source and its split
maps}\label{gammad-the-original-arithmetic-source-and-its-split-maps}

Retain

\[
G(R)=\{\tau_R\}\sqcup\{r^\bullet:r\in R\},\qquad e_R=0_R^\bullet,
\]

with the original maps

\[
\begin{array}{ccc}
G(\mathbb Z)&\xrightarrow{G(j)}&G(\mathbb C)\\
p_{\mathbb Z}\downarrow&&\downarrow p_{\mathbb C}\\
\mathbb Z&\xrightarrow{j}&\mathbb C.
\end{array}                                                    \tag{GD.1}
\]

The target of the first quotient is infinite. Its structural base is the
same absolute pointed base \(\mathfrak b_\tau\). The supported scalar
and the external absorber have not changed.

The source complex is

\[
C_+=[V\xrightarrow{\Theta}\mathscr B],\quad
Q=\mathscr B/\Theta V,\quad D=-x\partial_x,
\]

where \(V\) consists of the original even Schwartz functions with both
original moments zero, and \(\mathscr B\) has the original rapid decay
at both multiplicative ends with all \(D\)-derivatives. Keep

\[
\Theta\phi(x)=\sum_{n\ne0}\phi(nx),\qquad
 g(s)=2\xi(s)=s(s-1)\pi^{-s/2}\Gamma(s/2)\zeta(s).
\]

Fix a nonempty finite packet \(h(s)=\prod_\rho(s-\rho)^{m_\rho}\) of
actual zeros, with every selected order complete. For reflection
statements require the same dagger stability as the source. Write

\[
v_h=g/h,\quad \mathcal MF_h=v_h,\quad
w_h(t)=\frac{|v_h(1/2+it)|^2}{2\pi},\quad \mu_h=\int w_h(t)\,dt.       \tag{GD.2}
\]

The cyclic tensor source and its maps remain

\[
\mathcal V_{h,k}P=P(D_1+\cdots+D_k)(F_h^{\otimes k}),\quad
C_{h,k}=\mathbb C[S]/\chi_{h,k},\quad A=M_S,
\]

\[
J^{(k)}\mathcal V_{h,k}=\eta_{h,k}\pi_\chi,\qquad
q^{(k)}\mathcal V_{h,k}=\sigma_h^{\otimes k}\eta_{h,k}\pi_\chi,          \tag{GD.3}
\]

where \(S=s_1+\cdots+s_k\) and
\(\eta_{h,k}[P]=\upsilon_h^{\otimes k}P(A_k)1\), with
\(\upsilon_h=j_h(g/h)\). All local units remain in this map.

At polynomial degree \(N\), the original relation map is

\[
M_\chi:\mathcal P_{N-q}\longrightarrow\mathcal P_N,\quad P\mapsto\chi P,
\qquad q=\deg\chi,                                                   \tag{GD.4}
\]

and its image under \(\mathcal V_{h,k}\) consists of the already
constructed theta boundaries. Every norm construction below concerns
this same map and its quotient. Its squared-norm weight \(|\chi|^2\)
does not mean that the next conormal ideal is \((\chi^2)\): that ideal
remains the exact pullback of the original \(I^r\).

For \(h=1\), the analytic source (GD.2) still exists, but the finite
arithmetic quotient is the zero module. Its split lift has \(e\) and
\(\tau\). No nonzero packet or Gram inverse is inferred from that
analytic branch.

\subsection[\hspace{0.4em}Carry the repository's gamma reference through every tensor
power]{Carry the repository's gamma reference through every tensor
power}\label{gammad-carry-the-repositorys-gamma-reference-through-every-tensor-power}

For a fixed real \(\lambda>0\) define exactly

\[
r_\lambda(t)=\frac{|\Gamma(\lambda+it/2)|^2}{2\pi},\qquad
 d\sigma_\lambda(t)=r_\lambda(t)dt,\qquad
 c_\lambda=2^{1-2\lambda}\Gamma(2\lambda).                           \tag{GD.5}
\]

The repository proves, by the beta integral and Fourier inversion,

\[
\int e^{iyt}d\sigma_\lambda(t)=c_\lambda(\cosh y)^{-2\lambda}.
                                                                    \tag{GD.6}
\]

Its factor \(2\) comes from \(t=2u\) in the beta integral. Thus
\(c_\lambda\), rather than \(1\), is the mass. Fourier uniqueness for
finite measures gives the all-tensor identity

\[
\boxed{
 r_{\lambda,k}:=r_\lambda^{*k}
 =\kappa_{\lambda,k}r_{k\lambda},\qquad
 \kappa_{\lambda,k}=\frac{c_\lambda^k}{c_{k\lambda}}
 =2^{k-1}\frac{\Gamma(2\lambda)^k}{\Gamma(2k\lambda)}.
}                                                                   \tag{GD.7}
\]

In particular the mass of the sum measure is \(c_\lambda^k\). For
\(\lambda=1/4\) this reads

\[
\boxed{
 r_{1/4,k}(u)=
 \frac{2^{k-1}\pi^{k/2}}{\Gamma(k/2)}
 \frac{|\Gamma(k/4+iu/2)|^2}{2\pi},\qquad
 \int r_{1/4,k}=(2\pi)^{k/2}.
}                                                                   \tag{GD.8}
\]

The literal coordinate map is \((t_1,\ldots,t_k)\mapsto u=\sum_i t_i\);
\(u\) is not replaced by an average or a rescaled variable. The original
arithmetic scaling coordinate is still \(S=k/2+iu\).

For real \(|\theta|<\pi/2\), analytic continuation of (GD.6) gives

\[
Z_{\lambda,k}^{\Gamma}(\theta)
 :=\int e^{\theta u}r_{\lambda,k}(u)du
 =c_\lambda^k(\cos\theta)^{-2k\lambda}.                            \tag{GD.9}
\]

The exponential integrability follows either from the original
beta-contour estimate or the gamma vertical-strip bound. This justifies
all finite differentiated moment integrals used below.

\subsection[\hspace{0.4em}The sum projection and its complete polynomial
complement]{The sum projection and its complete polynomial
complement}\label{gammad-the-sum-projection-and-its-complete-polynomial-complement}

Define actual Hilbert spaces

\[
\mathcal H_k=L^2(\mathbb R^k,\prod_i d\sigma_\lambda(t_i)),\qquad
\mathcal H_\Sigma=L^2(\mathbb R,r_{\lambda,k}(u)du).
\]

Their isometric inclusion is

\[
U_\Sigma:\mathcal H_\Sigma\to\mathcal H_k,\qquad
(U_\Sigma f)(\mathbf t)=f(t_1+\cdots+t_k).                           \tag{GD.10}
\]

For \(k\ge2\), its adjoint has the explicit disintegration formula

\[
(C_\Sigma F)(u)=\frac{1}{r_{\lambda,k}(u)}
\int_{\mathbb R^{k-1}}
F(t_1,\ldots,t_{k-1},u-\sum_{i<k}t_i)
\prod_{i<k}r_\lambda(t_i)\,r_\lambda(u-\sum_{i<k}t_i)\,d^{k-1}t.
                                                                    \tag{GD.11}
\]

Its denominator is the actual fibre mass. Fubini proves
\(C_\Sigma U_\Sigma=1\) and \(C_\Sigma=U_\Sigma^*\). For \(k=1\) both
maps are the identity. Put \(P_\Sigma=U_\Sigma C_\Sigma\) and retain

\[
F\longmapsto(C_\Sigma F,(1-P_\Sigma)F),\qquad
(f,n)\longmapsto U_\Sigma f+n.                                      \tag{GD.12}
\]

These are inverse Hilbert isometries onto
\(\mathcal H_\Sigma\oplus\ker C_\Sigma\). The second component is not
discarded. At each original support label \(\ell\), their reconstructed
maps are \((\ell,F)\mapsto(\ell,C_\Sigma F)\) and
\((\ell,F)\mapsto(\ell,(1-P_\Sigma)F)\). A nonzero normal vector has
first image \((\ell,0)\), while its second image retains that vector.
Both maps send external absence to external absence; their paired
inverse is still (GD.12). Norm and Gram observations have their actual
quadratic types, not the type of a linear quotient:
\(\operatorname{Gram}:\operatorname{Hom}_{\mathbb C}(E,\mathcal H)
\to\operatorname{Herm}(E)\), \(R\mapsto R^*R\).

\subsubsection{Exact action on every product-polynomial
degree}\label{gammad-exact-action-on-every-product-polynomial-degree}

Let \(b_n^{(\lambda)}(t)\) be the repository's monic gamma polynomial,
with

\[
\sum_{n\ge0}\frac{b_n^{(\lambda)}(t)}{n!}z^n
=(1-iz)^{-\lambda+it/2}(1+iz)^{-\lambda-it/2}.
\]

Set \(\alpha=2\lambda\). Its complete norm and recurrence are

\[
\|b_n^{(\lambda)}\|^2=c_\lambda n!(\alpha)_n,\qquad
b_{n+1}^{(\lambda)}=t b_n^{(\lambda)}-n(n+\alpha-1)b_{n-1}^{(\lambda)}.
                                                                    \tag{GD.13}
\]

The polynomials are not divided by their norms. Multiplying the \(k\)
generating functions proves the literal addition formula

\[
b_n^{(k\lambda)}(\sum_i t_i)
=\sum_{n_1+\cdots+n_k=n}\frac{n!}{\prod_i n_i!}
                         \prod_i b_{n_i}^{(\lambda)}(t_i).         \tag{GD.14}
\]

For a multi-index \(\mathbf n\), put \(n=\sum_i n_i\). The exact sum
projection is

\[
\boxed{
C_\Sigma\left(\prod_i b_{n_i}^{(\lambda)}(t_i)\right)
 =\frac{\prod_i(\alpha)_{n_i}}{(k\alpha)_n}\,b_n^{(k\lambda)}(u).
}                                                                   \tag{GD.15}
\]

\textbf{Proof.} Pair the product against (GD.14), in the actual product
measure. Orthogonality makes the answer zero unless the sum degree is
\(n\), and at degree \(n\) the answer is
\(c_\lambda^k n!\prod_i(\alpha)_{n_i}\). The squared norm of the sum
polynomial is \(c_\lambda^k n!(k\alpha)_n\). Their quotient gives (GD.15).

These tests identify the whole conditional projection, not merely its
finite polynomial moments: the gamma sum measure has exponential
moments, so its polynomials are dense in \(L^2\). To verify that last
statement, an \(L^2\) vector orthogonal to all polynomials defines a
finite measure whose Fourier transform is analytic on a strip by
Cauchy--Schwarz and exponential integrability. All derivatives at zero
vanish; analytic uniqueness and uniqueness of Fourier transforms give
the zero vector.

The retained normal polynomial is

\[
n_{\mathbf n}=\prod_i b_{n_i}^{(\lambda)}(t_i)
 -\frac{\prod_i(\alpha)_{n_i}}{(k\alpha)_n}
                    b_n^{(k\lambda)}(\sum_i t_i).
\]

Its exact squared norm is

\[
\boxed{
\|n_{\mathbf n}\|^2
=c_\lambda^k\left[
\prod_i n_i!(\alpha)_{n_i}
 -\frac{n!\prod_i(\alpha)_{n_i}^2}{(k\alpha)_n}
\right].
}                                                                   \tag{GD.16}
\]

Thus every degree, every factorial, and the relative component have an
explicit map and norm. This is the classical Meixner--Pollaczek addition
mechanism, applied with the source's exact scale and masses; it is not a
claim to have discovered a new general family of orthogonal polynomials.

\subsection[\hspace{0.4em}Insert the complete arithmetic multiplier, not a substituted
gamma
measure]{Insert the complete arithmetic multiplier, not a substituted
gamma
measure}\label{gammad-insert-the-complete-arithmetic-multiplier-not-a-substituted-gamma-measure}

Define the complex amplitude and its squared modulus by

\[
A_{h,\lambda}(t)
=\frac{v_h(1/2+it)}{\Gamma(\lambda+it/2)},\qquad
B_{h,\lambda}(t)=|A_{h,\lambda}(t)|^2,
\qquad w_h=B_{h,\lambda}r_\lambda.                                 \tag{GD.17}
\]

The entire quotient \(g/h\) supplies values at the selected zeros. No
pole is inserted into the multiplier at those locations. The full
complex amplitude, not only its modulus, remains in the source
comparison

\[
P\longmapsto
\prod_i\frac{\Gamma(\lambda+it_i/2)}{\sqrt{2\pi}}
 A_{h,\lambda}(t_i)\,P(k/2+i\sum_i t_i).
                                                                    \tag{GD.18}
\]

This is the original Mellin image of \(\mathcal V_{h,k}P\).

The scalar multiplier on the sum line is

\[
B_{h,\lambda;k}=C_\Sigma(B_{h,\lambda}^{\otimes k}).
\]

The exact pushforward identity is

\[
\boxed{
 m_{h,k}(u)=w_h^{*k}(u)=r_{\lambda,k}(u)B_{h,\lambda;k}(u).
}                                                                   \tag{GD.19}
\]

This is a norm observation of the original source; it does not project
the original finite arithmetic module or change its trace. The maps in
(GD.3), including the full unit in \(\eta\), remain attached to every
polynomial. The full field \(B_{h,\lambda}^{\otimes k}\) is recoverable
from (GD.12), so its relative complement is retained too.

\subsubsection{A bounded multiplier for every full
quartet}\label{gammad-a-bounded-multiplier-for-every-full-quartet}

For \(\lambda=1/4\), the original formula is

\[
A_{h,1/4}(t)=
\frac{-\pi^{-1/4}(t^2+1/4)e^{-it\log\pi/2}\zeta(1/2+it)}{h(1/2+it)}.
                                                                    \tag{GD.20}
\]

The minus sign and phase are still in (GD.18). The elementary continuation
formula

\[
\zeta(s)=\frac{s}{s-1}-s\int_1^\infty\{x\}x^{-s-1}dx
\]

gives, on this line,

\[
|\zeta(1/2+it)|\le1+2\sqrt{t^2+1/4}\le2(1+|t|).
\]

Let \(d=\deg h\ge3\) and take \(T_h=\max(1,2\max_\rho|\rho-1/2|)\). For
\(|t|\ge T_h\), \(|h(1/2+it)|\ge(|t|/2)^d\), and direct substitution
gives

\[
B_{h,1/4}(t)\le\frac{25\,2^{2d}}{\sqrt\pi}|t|^{6-2d}.
\]

Consequently the finite, specified constant

\[
C_h=\max\left\{
\sup_{|t|\le T_h}|A_{h,1/4}(t)|^2,
\frac{25\,2^{2d}}{\sqrt\pi}T_h^{6-2d}
\right\}
\]

satisfies \(0\le B_{h,1/4}\le C_h\). The compact supremum uses the
analytic quotient, not the expression with a removable denominator left
untreated. A full nonreal quartet has \(d=4m\ge4\), so this applies to
the entire packet used in the exterior amplification, with constants
allowed to depend on that fixed packet.

The positive integral in (GD.11) now proves

\[
\boxed{
0<B_{h,1/4;k}(u)\le C_h^k\quad(k\ge2),\qquad
0\le m_{h,k}(u)\le C_h^k r_{1/4,k}(u).
}                                                                   \tag{GD.21}
\]

Strict positivity for \(k\ge2\) follows because the one-factor
multiplier only vanishes at a discrete set; its excluded hyperplanes
have zero fibre measure. For \(k=1\) retain its actual isolated zeros.

\subsubsection{The other packet sizes and the empty analytic
source}\label{gammad-the-other-packet-sizes-and-the-empty-analytic-source}

For a chosen reference below, \(C_h\) always denotes an explicitly
specified upper bound for that reference's \(B_{h,\lambda}\). An
explicit alternate reference is available without changing the
arithmetic source. For integer \(L\ge0\),

\[
A_{h,1/4+L}(t)
=\frac{A_{h,1/4}(t)}{\prod_{j=0}^{L-1}(j+1/4+it/2)}.                \tag{GD.22}
\]

Its product is a known nonvanishing gamma-shift factor on the
integration line. Taking \(L\ge\max(0,3-d)\) makes its squared modulus
bounded by the same compact/tail argument. Equations (GD.17)--(GD.19) connect
both references to the identical original \(w_h\); no coordinate or
original mass is reassigned.

For the actual seed \(h=1\), choosing \(L=3\) gives the explicit
all-real estimate

\[
\boxed{
B_{1,13/4}(t)\le1024/\sqrt\pi.
}                                                                   \tag{GD.23}
\]

Indeed the gamma-shift denominator is
\(4^{-3}(t^2+1/4)(t^2+25/4)(t^2+81/4)\), and
\(|\zeta(1/2+it)|\le4\sqrt{t^2+1/4}\). The remaining fraction is bounded
by one since \((t^2+1/4)^2\le(t^2+25/4)(t^2+81/4)\).

This is an analytic estimate for a nonzero seed. Its finite spectral
quotient is still zero, as required by the merged correction.

\subsection[\hspace{0.4em}Exact arithmetic coefficients at every tensor
degree]{Exact arithmetic coefficients at every tensor
degree}\label{gammad-exact-arithmetic-coefficients-at-every-tensor-degree}

Use a reference for which \(B=B_{h,\lambda}\) is bounded, as constructed
above. Define its actual coefficients

\[
c_{h,j}=\frac{\int b_j^{(\lambda)}(t)w_h(t)dt}
                 {c_\lambda j!(\alpha)_j},\qquad
\mathcal A_h(z)=\sum_{j\ge0}(\alpha)_j c_{h,j}z^j.                  \tag{GD.24}
\]

The series is used coefficientwise; no convergence radius for this
formal series is stipulated. Each coefficient is an actual one-factor
arithmetic integral. Define

\[
d_{h,k,n}=[z^n]\mathcal A_h(z)^k.
\]

Then the complete \(L^2(r_{\lambda,k}du)\) expansion is

\[
\boxed{
B_{h,\lambda;k}(u)
=\sum_{n\ge0}\frac{d_{h,k,n}}{(k\alpha)_n}
                                     b_n^{(k\lambda)}(u).
}                                                                   \tag{GD.25}
\]

\textbf{Proof.} Pair \(B^{\otimes k}\) with the finite addition formula
(GD.14). Each one-factor pairing is \(c_\lambda j!(\alpha)_j c_{h,j}\). The
multinomial coefficient cancels the displayed factorials and gives

\[
\int B_{h,\lambda;k} b_n^{(k\lambda)}r_{\lambda,k}
=c_\lambda^k n! [z^n]\mathcal A_h(z)^k.
\]

Divide by the exact norm \(c_\lambda^k n!(k\alpha)_n\). Completeness
gives (GD.25). No exchange of an infinite product of formal series with an
integral was needed: every coefficient used a finite polynomial
identity.

At degree zero,

\[
c_{h,0}=\mu_h/c_\lambda,\qquad
 d_{h,k,0}=(\mu_h/c_\lambda)^k,
\]

so (GD.25) recovers exactly \(\int m_{h,k}=\mu_h^k\).

The retained relative residual is

\[
\mathcal N_{h,k}=B^{\otimes k}-U_\Sigma B_{h,\lambda;k}.
\]

Its whole norm is calculated by Pythagoras:

\[
\boxed{
\|\mathcal N_{h,k}\|^2
=\left(\int B(t)^2d\sigma_\lambda(t)\right)^k
-c_\lambda^k\sum_{n\ge0}\frac{n!}{(k\alpha)_n}|d_{h,k,n}|^2.
}                                                                   \tag{GD.26}
\]

The difference is nonnegative because it is the norm of the explicitly
retained vector. It is not set to zero by calling the sum map a
pushforward.

\subsubsection{Finite exactness and approximation
errors}\label{gammad-finite-exactness-and-approximation-errors}

Let \(B_L=\sum_{j=0}^L c_{h,j}b_j^{(\lambda)}\) and
\(B_{k,L}=C_\Sigma(B_L^{\otimes k})\). Projection contraction and the
tensor telescoping identity give

\[
\boxed{
\|B_{h,\lambda;k}-B_{k,L}\|
\le k\|B\|^{k-1}\|B-B_L\|.
}                                                                   \tag{GD.27}
\]

These norms have the respective original reference measures. The proof
uses \(\|B_L\|\le\|B\|\) and expands the tensor difference into \(k\)
terms. Both sides retain the full mass factors.

More strongly, the moments of \(B_{k,L}r_{\lambda,k}\) against every
polynomial of degree at most \(L\) are \textbf{exactly} those of
\(m_{h,k}\), since coefficient \(n\le L\) in (GD.25) uses only \(c_{h,j}\)
with \(j\le n\).

It follows that the entire degree-\(N\) source Gram, and its
degree-\(N\) relation Gram, are determined exactly by
\(c_{h,0},\ldots,c_{h,2N}\). The relation integrands are also degree at
most \(2N\). Their tensor dependence is the explicit finite coefficient
power in (GD.24), not a fresh \(k\)-dimensional integral.

A truncated \(B_{k,L}\) need not be nonnegative pointwise. Its
relationship to the true positive metric is the exact finite-moment
equality just proved when \(L\ge2N\), or the error estimate (GD.27) and
Cauchy--Schwarz for the specified polynomial tests when \(L<2N\).
Positivity is not inferred from a plotted truncated density.

\subsubsection{Actual seed coefficients, and their tensor cross
term}\label{gammad-actual-seed-coefficients-and-their-tensor-cross-term}

The supplied numerical driver evaluates the original theta seed on
\(x\ge1\), using its inversion symmetry. With \(L=D-1/2\), the even
spectral moments are \(\mu_{2r}=2\int_1^\infty |L^rf_0(x)|^2dx\). It
constructs their source polynomials by \(q_0(z)=4z^2-6z\) and
\(q_{r+1}=2z(q_r-q_r')-q_r/2\), and integrates the integer double sums
using upper incomplete gamma integrals. This is a numerical check, not a
validated interval computation.

For \(h=1\) and the bounded-multiplier reference \(\lambda=13/4\), the
50-digit, cutoff-6 and 70-digit, cutoff-8 calculations agree in the
displayed values:

\begin{center}\begin{tabular}{@{}lr@{}}
\toprule
Quantity & Numerical value \\
\midrule
\(\mu_0\) & \(1.279007247846485140479533592267\) \\
\(\mu_2\) & \(13.055549302570558435392684658123\) \\
\(\mu_4\) & \(383.274341717364242554407237479131\) \\
\(\mu_6\) & \(17946.5273016052076235667646919834\) \\
\(c_{h,0}\) & \(0.201056688692785865809916163220\) \\
\(c_{h,2}\) & \(0.007645442999486401188133851481\) \\
\(c_{h,4}\) & \(-0.000030516920803762043746724840\) \\
\bottomrule\end{tabular}\end{center}

The reference mass is
\(c_{13/4}=6.361426004587219292629853548733\ldots\), so it has not been
assigned the actual mass \(\mu_0\) or the value one. The negative fourth
expansion coefficient is retained; a positive density need not have
positive coefficients in this orthogonal basis.

For every even arithmetic multiplier, the first tensor terms from (GD.25)
are exactly

\[
\begin{aligned}
d_{h,k,2}&=k c_{h,0}^{k-1}(\alpha)_2c_{h,2},\\
d_{h,k,4}&=k c_{h,0}^{k-1}(\alpha)_4c_{h,4}
 +\binom{k}{2}c_{h,0}^{k-2}\bigl((\alpha)_2c_{h,2}\bigr)^2.
\end{aligned}
\]

Both terms in the fourth coefficient remain. These are coefficients of
the exact arithmetic source norm; the empty packet still supplies no
nonzero finite spectral quotient. Agreement of two numerical truncations
is not a bound on their shared truncation or rounding error.

\subsection[\hspace{0.4em}Feed the result into the two original Toda
determinants]{Feed the result into the two original Toda
determinants}\label{gammad-feed-the-result-into-the-two-original-toda-determinants}

Put \(\alpha_k=2k\lambda\). The full gamma source determinants are
explicit:

\[
\boxed{
\mathfrak D_n^\Gamma(\theta)
=c_\lambda^{kn}
\prod_{j=0}^{n-1}j!(\alpha_k)_j
(\cos\theta)^{-n(\alpha_k+n-1)}.
}                                                                   \tag{GD.28}
\]

At \(\theta=0\) this is the product of the monic norms from (GD.13) and
(GD.7). One direct proof for all real \(\theta\) starts with
\(\mathfrak D_0^\Gamma=1\) and
\(\mathfrak D_1^\Gamma=c_\lambda^k\cos^{-\alpha_k}\theta\). Apply the
already-proved Hankel--Toda identity recursively; the second derivative
of the logarithm of the right side is \(n(\alpha_k+n-1)\sec^2\theta\).
Substitution proves the next determinant with all constants. Strict
positivity makes the recurrence division valid.

In particular

\[
\omega_n^\Gamma(\theta)
=c_\lambda^k n!(\alpha_k)_n(\cos\theta)^{-(\alpha_k+2n)},\qquad
 a_n^\Gamma(\theta)=n(n+\alpha_k-1)\sec^2\theta.                     \tag{GD.29}
\]

The arithmetic source determinant is the previous \(\mathfrak D_n\). The
reference relation determinant is

\[
\mathfrak B_n^\Gamma(\theta)
=\det\left[\partial_\theta^{i+j}
 \left\{\overline\chi(k/2-i\partial_\theta)
       \chi(k/2+i\partial_\theta)
       Z_{\lambda,k}^\Gamma(\theta)\right\}\right]_{i,j<n}.         \tag{GD.30}
\]

It is the Gram of the \textbf{same} multiplication map \(M_\chi\). This
is an explicit finite expression in gamma constants, \(\cos\theta\),
\(\tan\theta\), and the original coefficients of \(\chi\).

Define positive scalar corrections, including the empty-index value one,

\[
X_n=\mathfrak D_n/\mathfrak D_n^\Gamma,\qquad
Y_n=\mathfrak B_n/\mathfrak B_n^\Gamma.
\]

The source--boundary determinant theorem now factors as

\[
\boxed{
V_N=V_N^\Gamma\,T_N,\qquad
 V_N^\Gamma=\frac{\mathfrak D_{N+1}^\Gamma}{\mathfrak B_{N-q+1}^\Gamma},\qquad
 T_N=\frac{X_{N+1}}{Y_{N-q+1}}.
}                                                                   \tag{GD.31}
\]

This retains the entire arithmetic correction in both numerator and
denominator. It is not an assertion that the arithmetic quotient metric
equals the reference metric.

\subsubsection{The representative map between the two
metrics}\label{gammad-the-representative-map-between-the-two-metrics}

Use the same finite polynomial source. Let \(B_\chi\) be the injective
matrix of its original relation columns, and let \(R_N^\Gamma\) be the
least-norm lift for the gamma observation. Let \(M_h\) be the actual
arithmetic source Gram. Then the actual lift is

\[
\boxed{
R_N=R_N^\Gamma-
B_\chi(B_\chi^*M_hB_\chi)^{-1}B_\chi^*M_hR_N^\Gamma.
}                                                                   \tag{GD.32}
\]

For an empty boundary block the correction is the zero map. The proof is
direct: the correction is an old relation, so the quotient remains the
identity; multiplication by \(B_\chi^*M_h\) gives zero. These two
properties characterize the arithmetic least-norm lift.

Applying \(\mathcal V_{h,k}\) to (GD.32) supplies the original theta
primitive of the difference. In the reconstructed support carrier, the
difference maps to the supported zero in its fibre. It does not
disappear from the source metric or from its boundary primitive. The
comparison keeps the complete unit in (GD.3).

\textbf{The boundary and its explicit cochain primitive.}
For a class \(x\), let
\[
P_x=-(B_\chi^*M_hB_\chi)^{-1}B_\chi^*M_hR_N^\Gamma x
\]
denote its coefficient vector in \(\mathcal P_{N-q}\).
Then \((R_N-R_N^\Gamma)x=\chi P_x\).
In the preceding sentence, applying \(\mathcal V_{h,k}\) gives the
\emph{boundary}; its primitive is obtained as follows, retaining its
cochain degree. Use the actual successive monic divisions
\(\chi(S)=\sum_i h(s_i)Q_i(\mathbf s)\) of TVB.8 and put
\[
\begin{split}
\mathfrak p_x={}&\sum_{i=1}^k(-1)^{i-1}
 Q_i(D_1,\ldots,D_k)P_x(D_1+\cdots+D_k)\\
&\hspace{0.4in}\cdot
 \bigl(F_h^{\otimes(i-1)}\otimes\phi_0
                  \otimes F_h^{\otimes(k-i)}\bigr),\\
\phi_0(x)={}&(4\pi^2x^4-6\pi x^2)e^{-\pi x^2}.
\end{split} \tag{GD.32a}
\]
Here \(\mathfrak p_x\) has tensor cochain degree \(k-1\).
Since \(h(D)F_h=\Theta\phi_0\), the tensor differential contributes
the second factor \((-1)^{i-1}\) and gives
\[
d\mathfrak p_x=\mathcal V_{h,k}(\chi P_x)
              =\mathcal V_{h,k}(R_N-R_N^\Gamma)x. \tag{GD.32b}
\]
This proves the asserted relation without identifying a boundary with its
primitive. When the relation block is empty, both maps are zero.


\subsubsection{A genuine source-side upper
bound}\label{gammad-a-genuine-source-side-upper-bound}

From (GD.21), for every admitted polynomial,

\[
\|P\|_h^2\le C_h^k\|P\|_\Gamma^2.
\]

Evaluate this inequality on \(R_N^\Gamma u\) and then minimize on the
arithmetic side. This proves

\[
\boxed{
G_N\preceq C_h^kG_N^\Gamma,\qquad
0<T_N\le C_h^{kq}.
}                                                                   \tag{GD.33}
\]

The second assertion is the determinant inequality in dimension \(q\).
It has the quotient dimension, rather than the full polynomial-source
dimension, in its exponent. The corresponding source estimate is
\(0<X_n\le C_h^{kn}\), and similarly \(0<Y_n\le C_h^{kn}\).

These are one-sided estimates. They do not bound the consecutive ratio
\(T_{N-1}/T_{N+1}\), which is the quantity required by the Toda upper
bound. Its exact map to that estimate is the following.

\subsection[\hspace{0.4em}The current upper-control expression, with all reference
costs
exposed]{The current upper-control expression, with all reference
costs
exposed}\label{gammad-the-current-upper-control-expression-with-all-reference-costs-exposed}

Define

\[
Q_N=\frac{X_{N+2}X_N}{X_{N+1}^2},\qquad
\mathcal R_N^\Gamma=\frac{V_{N-1}^\Gamma}{V_{N+1}^\Gamma}.
\]

At the original observation \(\theta=0\), (GD.29)--(GD.31) give

\[
\boxed{
a_{N+1}=(N+1)(N+\alpha_k)Q_N,\qquad
\mathcal R_N=\mathcal R_N^\Gamma\frac{T_{N-1}}{T_{N+1}}.
}                                                                   \tag{GD.34}
\]

The preceding exact phase identity is also retained through

\[
\ell_N'=\partial_\theta\log V_N^\Gamma+\partial_\theta\log T_N.
\]

Thus its finite upper bound becomes exactly

\[
\boxed{
\epsilon_{h,k,N}\le
\sqrt{(N+1)(N+\alpha_k)Q_N}\,
\frac{\mathcal R_N^\Gamma T_{N-1}/T_{N+1}-1}
 {2\sqrt{\mathcal R_N^\Gamma T_{N-1}/T_{N+1}}}.
}                                                                   \tag{GD.35}
\]

The nonnegative numerator uses the original monotonicity of the actual
quotient minimum, not a presumed monotonicity of \(T_N\). When this
inequality is coarse, the inherited equality with the volume-imbalance
square and the phase square remains the sharper calculation.

For the full quartet one can keep \(\lambda=1/4\), so \(\alpha_k=k/2\).
The gamma source recurrence is now explicit at every degree; the
remaining arithmetic dependence is in \(Q_N\) and the relative
source--boundary correction \(T_{N-1}/T_{N+1}\). The reference relation
term still contains the full spectral annihilator \(\chi\). It is not
bounded by pretending its roots are on the integration line.

The exterior theorem remains

\[
L_{h,k}\le\epsilon_{h,k,N},\qquad
L_{h,k}=2\delta[1+k(m-1)](k+1)
          \left\lfloor\frac{(k+1)^2}{4}\right\rfloor.
\]

Neither (GD.33) nor (GD.35) proves a subcubic bound. Equation (GD.35) identifies
exactly which degree-to-degree arithmetic corrections must be controlled
next. All entries of those finite determinants can now be generated by
the all-\(k\) coefficient transform (GD.25), with the source--boundary
comparison map (GD.32) retained.

\subsection[\hspace{0.4em}A proved all-degree tail estimate for the actual arithmetic
moments]{A proved all-degree tail estimate for the actual arithmetic
moments}\label{gammad-a-proved-all-degree-tail-estimate-for-the-actual-arithmetic-moments}

The multiplier bound gives more than a change of coordinates. For
\(a,b>0\) with \(a+b<\pi/2\), every integer \(r\ge0\) and \(T\ge0\),

\[
\boxed{
\int_{|u|>T}|u|^r m_{h,k}(u)du
\le
2 C_h^k c_\lambda^k\,r!\,a^{-r}
 e^{-bT}(\cos(a+b))^{-2k\lambda}.
}                                                                   \tag{GD.36}
\]

\textbf{Proof.} For \(x\ge0\), \(x^r\le r!a^{-r}e^{ax}\). On the
indicated tail, \(1\le e^{b(x-T)}\). Apply (GD.21) and integrate. Since the
reference density is even,

\[
\int e^{(a+b)|u|}r_{\lambda,k}(u)du
\le Z_{\lambda,k}^\Gamma(a+b)+Z_{\lambda,k}^\Gamma(-a-b).
\]

Equation (GD.9) gives (GD.36). This estimate retains the original mass and the
complete tensor cost. It is a theorem with specified constants, not a
claim that a floating-point quadrature has validated them numerically.

For a source or relation moment whose integrand is a polynomial
\(P(u)=\sum_j p_ju^j\), sum \(|p_j|\) times (GD.36). In particular this
controls every entry of a degree-\(N\) source and relation Gram,
including the full \(|\chi|^2\) factor. It is then combined with an
independently certified compact-interval quadrature and the existing
matrix inverse/error inequalities; it does not certify that quadrature
by itself.

This is an explicit tail interface for the actual arithmetic input at
all degrees and tensor powers. For the empty analytic seed, (GD.23) gives a
completely stated bound without locating any zeta zero. For a nonempty
packet the compact constant in \(C_h\) must be enclosed using the actual
analytic quotient.

\subsection[\hspace{0.4em}What has and has not been
concluded]{What has and has not been
concluded}\label{gammad-what-has-and-has-not-been-concluded}

The constructed chain is

\[
\begin{gathered}
\text{original split-supported polynomial source}\\
\longrightarrow\text{full gamma product fibre with arithmetic amplitude}\\
\longrightarrow\text{sum density and retained relative component}\\
\longrightarrow\text{exact all-tensor arithmetic coefficients}
\\
\longrightarrow\text{the same source--boundary Toda determinant ratio}.
\end{gathered}
\]

The finite coefficient projection, polynomial degree, all masses, all
gamma factors, the arithmetic unit, and the original theta primitives
are explicit. The same finite quotient is used throughout; an existing
relation becomes \(e\) only under its original quotient map, while the
source norm and all derivative data remain available before that map. No
new scalar-zero adjunction is performed at a tensor or relation level.

There is now a proved bound for the full arithmetic sum density,
explicit closed-form reference determinants at every tensor degree, an
exact finite coefficient algorithm, and a rigorous moment-tail bound.
There is not yet a uniform small bound for the consecutive arithmetic
determinant correction. The upper-estimate problem is sharpened to that
specific comparison rather than being claimed solved by gamma
convolution or by positivity of individual source Grams.

\subsection[\hspace{0.4em}References and exact source
roles]{References and exact source
roles}\label{gammad-references-and-exact-source-roles}

\begin{itemize}
\tightlist
\item
  Repository main at \path{fa4be32f87a9a90f3c02a07f7790dea95a1e7b0f},
  \path{tex/theta_gamma_reference.tex} in the 20260912-stieltjes
  continuation, TG.2--TG.14: original one-factor gamma reference,
  arithmetic amplitude, constants and polynomial norms. Only the stated
  TeX window was read in this review; no complete reader audit is
  claimed.
\item
  PR \#22 head \path{811210d24b80813a08972ca23f919db015383137},\par
  \path{workbenches/tau-exterior-trace-formal/RESEARCH_NOTE.md}:
  current exterior and filtration formalization work, with its separate
  verification scope. Its reported in-progress status is not promoted to
  a completed certificate.
\item
  Delivered \path{Tau_Toda_Volume_Control/NOTE.tex}: original
  source and relation determinants, quotient metric, phase derivative
  and finite upper inequality. Its supplied manifest was independently
  checked here.
\item
  NIST Digital Library of Mathematical Functions, \S\S18.22 and 18.23:
  classical Meixner--Pollaczek recurrence and generating function. The
  formulas here retain the repository's original \(t/2\) scale and all
  factorials.
\item
  Koelink, H. T., and Van der Jeugt, J. (1998). Convolutions for
  orthogonal polynomials from Lie and quantum algebra representations.
  \emph{SIAM Journal on Mathematical Analysis}, 29(3), 794--822.
  arXiv:q-alg/9607010. Background for the classical addition/convolution
  mechanism, not a claim that this article treats the present arithmetic
  quotient.
\end{itemize}

The proofs in this note establish the stated specialized maps directly.
The regression checker tests finite exact polynomial and matrix
consequences; it does not replace the analytic proofs or establish the
outstanding uniform estimate.


\endgroup

\subsection[\hspace{0.4em}The analytic coefficient map into the same Toda input]{The analytic coefficient map into the same Toda input}
\label{sec:gamma-toda-analytic-coefficient-map}

The preceding coefficient construction initially treats its power series
formally. The exponential decay of the actual theta source proves more:
the series is an analytic re-expression of the same arithmetic moment
function, with an explicit inverse near the original observation.
This joins the classical Meixner--Pollaczek generating function to the
arithmetic source. It does not replace the arithmetic measure by the
reference measure.

Keep
\[
\begin{split}
M_h(\theta)&=\int_{\mathbb R}e^{\theta t}w_h(t)\,dt,
\qquad w_h(t)=\frac{|(g/h)(1/2+it)|^2}{2\pi},\\
c_\lambda&=2^{1-2\lambda}\Gamma(2\lambda),\qquad
\alpha=2\lambda,\qquad
\mathcal A_h(z)=\sum_{n\geq0}(\alpha)_n c_{h,n}z^n.
\end{split}
\]
Here \(c_{h,n}\) are exactly GD.24, with the original monic polynomials
\(b_n^{(\lambda)}\) and their actual norms. In particular
\(\mathcal A_h(0)=\mu_h/c_\lambda\), not one in general.

\begin{proposition}[Exact analytic generating-function map]
Use logarithms analytic on \(|z|<1\) and vanishing at \(z=0\) for
\(1+iz\) and \(1-iz\). Then the coefficient series is holomorphic on
that disk and
\[
\boxed{\mathcal A_h(z)=
 \frac{(1+z^2)^{-\lambda}}{c_\lambda}
 M_h(\arctan z).} \tag{GJ.1}
\]
For real \(|\theta|<\pi/4\) its inverse is
\[
\boxed{M_h(\theta)=c_\lambda(\sec\theta)^{2\lambda}
                 \mathcal A_h(\tan\theta).} \tag{GJ.2}
\]
Consequently, at every integer \(k\geq1\), with the same unscaled sum
\(u=t_1+\cdots+t_k\),
\[
\boxed{\frac{Z_{h,k}(\theta)}{Z^\Gamma_{\lambda,k}(\theta)}
       =\mathcal A_h(\tan\theta)^k,\quad
 Z^\Gamma_{\lambda,k}(\theta)
       =c_\lambda^k(\sec\theta)^{2k\lambda}.} \tag{GJ.3}
\]
The last identity is asserted on the stated real interval, and as an
identity of holomorphic germs at zero.
\end{proposition}

\begin{proof}
The classical generating function, in the coordinates of GD.13, is
\[
\sum_{n\geq0}\frac{b_n^{(\lambda)}(t)}{n!}z^n
 =(1-iz)^{-\lambda+it/2}(1+iz)^{-\lambda-it/2}
 =(1+z^2)^{-\lambda}e^{t\arctan z}. \tag{GJ.4}
\]
This is DLMF 18.23.7 at \(\phi=\pi/2\), \(x=t/2\), with
\(b_n=n!P_n^{(\lambda)}\); the factors of two are therefore retained.
For \(|z|<1\),
\[
W(z)=\frac{1+iz}{1-iz},\qquad
\operatorname{Re}W(z)=\frac{1-|z|^2}{|1-iz|^2}>0,
\qquad \arctan z=\frac{\operatorname{Log}W(z)}{2i}.
\]
Thus \(|\operatorname{Re}\arctan z|<\pi/4\).
Also \(\operatorname{Log}(1+z^2)=
\operatorname{Log}(1-iz)+\operatorname{Log}(1+iz)\) on this disk:
the equality holds at zero and its two analytic sides have the same
derivative.

The original arithmetic input \(M_h\) is holomorphic for
\(|\operatorname{Re}\theta|<\pi/2\), by the complete zero division and
vertical-strip gamma bound proved in the Toda calculation.
On a compact subset of the \(z\)-disk, the absolute value of the integrand
in GJ.4 times \(w_h(t)\) is bounded by a fixed constant times
\(e^{v|t|}w_h(t)\), with \(v<\pi/2\).
The same domination, with a polynomial factor in \(|t|\), applies to
each fixed derivative. It justifies holomorphic integration and
differentiation at zero. The Taylor coefficient of
\[
\frac1{c_\lambda}\int_{\mathbb R}
 (1+z^2)^{-\lambda}e^{t\arctan z}w_h(t)\,dt
\]
is exactly
\[
\frac1{c_\lambda n!}\int b_n^{(\lambda)}(t)w_h(t)\,dt
       =(\alpha)_n c_{h,n}.
\]
This proves GJ.1, including convergence, without an unsupported exchange
of an infinite series and an integral. For real \(|\theta|<\pi/4\),
\(z=\tan\theta\) belongs to the disk,
\(\arctan(\tan\theta)=\theta\), and
\(1+\tan^2\theta=\sec^2\theta\), with positive real powers.
This proves GJ.2. Taking the literal \(k\)-th power and using GD.9 gives
GJ.3. At zero it gives \(Z_{h,k}(0)=\mu_h^k\), as required.
\end{proof}

\begin{proposition}[Finite coefficient jets for the actual metric]
Fix the original monic relation \(\chi\) of degree \(q\geq1\), an
admitted degree \(N\geq q-1\), and a derivative order \(j\geq0\).
The source and relation Gram matrices, and all their derivatives through
order \(j\) at \(\theta=0\), are determined exactly by
\[
\boxed{c_{h,0},\ldots,c_{h,2N+j},\quad k,\lambda,
        \text{and the original coefficients of }\chi.} \tag{GJ.5}
\]
The same finite dependence holds for the canonical quotient metric,
its determinant and its derivatives through order \(j\).
For \(N\geq q\), the adjacent-volume control in GD.35 at zero requires
only \(c_{h,0},\ldots,c_{h,2N+2}\), together with the fixed algebraic
packet data. This is a finite-dependence assertion, not a numerical
conditioning bound.
\end{proposition}

\begin{proof}
Write \(c=k/2\). In original \(S\)-coordinates, a source Gram entry
and its \(j\)-th derivative are
\[
\left.\partial_\theta^j H_{ab}(\theta)\right|_{\theta=0}
 =\int_{\mathbb R}(c-iu)^a(c+iu)^b u^j m_{h,k}(u)\,du,
 \qquad 0\leq a,b\leq N. \tag{GJ.6}
\]
The integrand polynomial has degree at most \(2N+j\).
A relation Gram entry replaces the two factors by
\(\overline{\chi(c+iu)(c+iu)^a}\) and
\(\chi(c+iu)(c+iu)^b\), where \(0\leq a,b\leq N-q\).
Its degree is at most \(2q+2(N-q)+j=2N+j\); the original two conjugate
factors and both signs remain.

By GD.25, the pairing of a polynomial of degree \(d\) with the actual
sum density uses exactly the expansion coefficients through degree
\(d\). Each such coefficient is
\([z^n]\mathcal A_h(z)^k/(k\alpha)_n\), and the coefficient of degree
\(n\) of the finite product uses only \(c_{h,0},\ldots,c_{h,n}\).
Equivalently, GJ.3 computes the same finite moment jet through composition
with \(\tan\theta\), whose constant term is zero.

In the ordered source basis
\(1,S,\ldots,S^{q-1},\chi,\chi S,\ldots,\chi S^{N-q}\),
write the source Gram as
\[
H(\theta)=\begin{pmatrix}H_{00}&H_{01}\\H_{10}&H_{11}\end{pmatrix},
\qquad
G_N(\theta)=H_{00}-H_{01}H_{11}^{-1}H_{10}. \tag{GJ.7}
\]
For \(N\geq q\), \(H_{11}(0)\) is positive definite, being the Gram
of linearly independent original relations against the positive
arithmetic measure. The inverse derivative is computed without new
moments by differentiating \(H_{11}H_{11}^{-1}=I\). Explicitly, if
\(K=H_{11}^{-1}\),
\[
K^{(r)}(0)=-K(0)\sum_{s=1}^r\binom rs
 H_{11}^{(s)}(0)K^{(r-s)}(0),\qquad r\geq1. \tag{GJ.8}
\]
Thus the Schur complement and its determinant have the asserted finite
dependence. At \(N=q-1\) the relation block is empty and \(G_N=H_{00}\),
so no inverse is required. Finally, \(a_{N+1}\) and \(V_{N+1}\) require
source degree \(N+1\), hence moments through \(2N+2\); the phase term
\(\partial_\theta\log V_N\) uses only degree \(2N+1\).
All lower degrees are already included. This proves the last assertion.
\end{proof}

\textbf{Consequence for the calculation.}
The all-tensor coefficient formula and the Toda determinant derivatives
are two explicit descriptions of the same one-factor arithmetic input.
At any fixed degree this gives a finite route from one-factor coefficients
to the original quotient metric and control, including their derivative
terms. The number of required coefficients still grows with the admitted
degree; no estimate uniform in growing \(N\) and \(k\) follows merely
from this finite identity.

\textbf{Source and verification scope.}
The generating function is classical
(\href{https://dlmf.nist.gov/18.23.E7}{DLMF 18.23.7}).
The integral map, its domain, inverse, and finite-jet consequence above
were calculated here from that identity and the actual source in GD.2.
They were independently checked with all mass and coordinate factors.
The package's 17 normal and 17 optimized finite tests are regression
checks of its delivered formulas, not Lean or interval certificates of
these analytic arguments or of a growing-degree upper estimate.
